\documentclass[10pt,a4paper]{amsart}
\usepackage{amsmath,amssymb,amsthm,graphicx}
% Retained minimal preamble for the Stacks Project chapter flat.tex (More on Flatness).
% Verbatim excerpt of the chapter preamble: Xy-pic setup (native preamble lines 14-18) and the native
% theorem declarations with the mathematical macros (native preamble lines 154-217). Copied mechanically;
% nothing added, removed or redefined. The wrapper class and the amsmath/amssymb/amsthm/graphicx base
% packages are supplied by the assembler, so the class/externaldocument part of the native preamble is omitted.
% For commutative diagrams we use Xy-pic
\usepackage[all]{xy}

% We use 2cell for 2-commutative diagrams.
\xyoption{2cell}
% Theorem environments.
%
\theoremstyle{plain}
\newtheorem{theorem}[subsection]{Theorem}
\newtheorem{proposition}[subsection]{Proposition}
\newtheorem{lemma}[subsection]{Lemma}

\theoremstyle{definition}
\newtheorem{definition}[subsection]{Definition}
\newtheorem{example}[subsection]{Example}
\newtheorem{exercise}[subsection]{Exercise}
\newtheorem{situation}[subsection]{Situation}

\theoremstyle{remark}
\newtheorem{remark}[subsection]{Remark}
\newtheorem{remarks}[subsection]{Remarks}

\numberwithin{equation}{subsection}

% Macros
%
\def\lim{\mathop{\mathrm{lim}}\nolimits}
\def\colim{\mathop{\mathrm{colim}}\nolimits}
\def\Spec{\mathop{\mathrm{Spec}}}
\def\Hom{\mathop{\mathrm{Hom}}\nolimits}
\def\Ext{\mathop{\mathrm{Ext}}\nolimits}
\def\SheafHom{\mathop{\mathcal{H}\!\mathit{om}}\nolimits}
\def\SheafExt{\mathop{\mathcal{E}\!\mathit{xt}}\nolimits}
\def\Sch{\mathit{Sch}}
\def\Mor{\mathop{\mathrm{Mor}}\nolimits}
\def\Ob{\mathop{\mathrm{Ob}}\nolimits}
\def\Sh{\mathop{\mathit{Sh}}\nolimits}
\def\NL{\mathop{N\!L}\nolimits}
\def\CH{\mathop{\mathrm{CH}}\nolimits}
\def\proetale{{pro\text{-}\acute{e}tale}}
\def\etale{{\acute{e}tale}}
\def\QCoh{\mathit{QCoh}}
\def\Ker{\mathop{\mathrm{Ker}}}
\def\Im{\mathop{\mathrm{Im}}}
\def\Coker{\mathop{\mathrm{Coker}}}
\def\Coim{\mathop{\mathrm{Coim}}}

% Boxtimes
%
\DeclareMathSymbol{\boxtimes}{\mathbin}{AMSa}{"02}

%
% Macros for moduli stacks/spaces
%
\def\QCohstack{\mathcal{QC}\!\mathit{oh}}
\def\Cohstack{\mathcal{C}\!\mathit{oh}}
\def\Spacesstack{\mathcal{S}\!\mathit{paces}}
\def\Quotfunctor{\mathrm{Quot}}
\def\Hilbfunctor{\mathrm{Hilb}}
\def\Curvesstack{\mathcal{C}\!\mathit{urves}}
\def\Polarizedstack{\mathcal{P}\!\mathit{olarized}}
\def\Complexesstack{\mathcal{C}\!\mathit{omplexes}}
% \Pic is the operator that assigns to X its picard group, usage \Pic(X)
% \Picardstack_{X/B} denotes the Picard stack of X over B
% \Picardfunctor_{X/B} denotes the Picard functor of X over B
\def\Pic{\mathop{\mathrm{Pic}}\nolimits}
\def\Picardstack{\mathcal{P}\!\mathit{ic}}
\def\Picardfunctor{\mathrm{Pic}}
\def\Deformationcategory{\mathcal{D}\!\mathit{ef}}
\usepackage{cleveref}
\crefname{originallemma}{original lemma}{original lemmas}
\Crefname{originallemma}{Original lemma}{Original lemmas}
\makeatletter
\newlabel{lemma-homothety-spectrum}{{05DE}{source}{}{sourcepointer-1}{}}
\newlabel{lemma-homothety-spectrum@cref}{{[originallemma][0][]05DE}{[1][0][]source}{}{}{}}
\makeatother
\crefname{originallemma}{original lemma}{original lemmas}
\Crefname{originallemma}{Original lemma}{Original lemmas}
\makeatletter
\newlabel{lemma-homothety-universally-injective}{{05DF}{source}{}{sourcepointer-2}{}}
\newlabel{lemma-homothety-universally-injective@cref}{{[originallemma][0][]05DF}{[1][0][]source}{}{}{}}
\makeatother
\crefname{originallemma}{original lemma}{original lemmas}
\Crefname{originallemma}{Original lemma}{Original lemmas}
\makeatletter
\newlabel{algebra-lemma-universally-injective-permanence}{{05CJ}{source}{}{sourcepointer-3}{}}
\newlabel{algebra-lemma-universally-injective-permanence@cref}{{[originallemma][0][]05CJ}{[1][0][]source}{}{}{}}
\makeatother
\crefname{originallemma}{original lemma}{original lemmas}
\Crefname{originallemma}{Original lemma}{Original lemmas}
\makeatletter
\newlabel{algebra-lemma-ff-rings}{{00HQ}{source}{}{sourcepointer-4}{}}
\newlabel{algebra-lemma-ff-rings@cref}{{[originallemma][0][]00HQ}{[1][0][]source}{}{}{}}
\makeatother
\crefname{originallemma}{original lemma}{original lemmas}
\Crefname{originallemma}{Original lemma}{Original lemmas}
\makeatletter
\newlabel{algebra-lemma-faithfully-flat-universally-injective}{{05CK}{source}{}{sourcepointer-5}{}}
\newlabel{algebra-lemma-faithfully-flat-universally-injective@cref}{{[originallemma][0][]05CK}{[1][0][]source}{}{}{}}
\makeatother
\crefname{originallemma}{original lemma}{original lemmas}
\Crefname{originallemma}{Original lemma}{Original lemmas}
\makeatletter
\newlabel{algebra-lemma-universally-injective-tensor}{{05CH}{source}{}{sourcepointer-6}{}}
\newlabel{algebra-lemma-universally-injective-tensor@cref}{{[originallemma][0][]05CH}{[1][0][]source}{}{}{}}
\makeatother
\crefname{originallemma}{original lemma}{original lemmas}
\Crefname{originallemma}{Original lemma}{Original lemmas}
\makeatletter
\newlabel{algebra-lemma-composition-universally-injective}{{05CI}{source}{}{sourcepointer-7}{}}
\newlabel{algebra-lemma-composition-universally-injective@cref}{{[originallemma][0][]05CI}{[1][0][]source}{}{}{}}
\makeatother
\title{Universally injective base change of a flat finite presentation module over a local homomorphism}
\author{The Stacks Project Authors}
\date{stacks-project:flat.tex:local-snapshot-20261011; GNU FDL 1.2 or later}
\begin{document}
\maketitle
\noindent\textit{Editorial note.} All delivered mathematical text is the author's own, copied from the pinned Stacks Project chapter file flat.tex; no statement or proof step was rewritten, reordered or removed.\par\medskip
\noindent\textit{Editorial note.} The preamble is the chapter's own retained minimal preamble (sources/stacks-flat/retained-preamble.tex), a verbatim excerpt of the native Xy-pic setup and of the native theorem declarations and mathematical macros; the Xy-pic xy environment is required because the counted proof opens with a displayed commutative square. No macro was invented or redefined.\par\medskip
\noindent\textit{Editorial note.} Same-chapter referenced results are located at their \textbackslash{}label line inside the pinned chapter file; the five cross-chapter algebra lemmas are located in the project's own tag list staged in the snapshot (sources/stacks-chow/tags.txt), at the verbatim lines carrying their canonical tags.\par\medskip
\noindent\textit{Allowed standard prerequisites (editorial).}\begin{itemize}
\item The semi-local localisation attached to a flat finite presentation module. Same chapter, Lemma labelled lemma-homothety-spectrum, source line 1783 (Stacks tag 05DE). Supplies that Sigma\textasciicircum{}\{-1\}S is semi-local with maximal ideals in bijection with the associated primes of N/m\_R N; the faithful flatness conclusion drawn from the hypothesis on Spec(S') -> Spec(S) is the counted step.
\item Universal injectivity of the localisation map. Same chapter, Lemma labelled lemma-homothety-universally-injective, source line 1821 (Stacks tag 05DF). Supplies that N -> Sigma\textasciicircum{}\{-1\}N is R-universally injective; the composition argument that finishes the proof stays inside the counted span.
\item Standard permanence and tensor facts for universally injective maps and faithfully flat ring maps. Algebra, Lemmas algebra-lemma-universally-injective-permanence (05CJ), algebra-lemma-ff-rings (00HQ), algebra-lemma-faithfully-flat-universally-injective (05CK), algebra-lemma-universally-injective-tensor (05CH) and algebra-lemma-composition-universally-injective (05CI). Cited at proof lines 1877-1902 for permanence under a further localisation, for faithful flatness from the condition on spectra, for universally injective maps induced by faithfully flat ring maps and for composition; none of them contains the reduction or the diagram chase of this proof.
\end{itemize}\medskip
\ifdefined\hypertarget\hypertarget{sourcepointer-1}{}\fi
\noindent\textit{Source reference.} Original lemma 05DE: sources/stacks-flat/pinned-original-source.tex, source lines 1783--1783. Same-chapter pointer invoked at proof line 1882 for semi-locality of Sigma\textasciicircum{}\{-1\}S and the bijection with associated primes; the deduction of faithful flatness from the hypothesis on the image of Spec(S') -> Spec(S) is retained in the counted span.\par\medskip
\ifdefined\hypertarget\hypertarget{sourcepointer-2}{}\fi
\noindent\textit{Source reference.} Original lemma 05DF: sources/stacks-flat/pinned-original-source.tex, source lines 1821--1821. Same-chapter pointer invoked at proof line 1898 for universal injectivity of N -> Sigma\textasciicircum{}\{-1\}N; the composition step that turns this into the conclusion for N -> N' is retained in the counted span.\par\medskip
\ifdefined\hypertarget\hypertarget{sourcepointer-3}{}\fi
\noindent\textit{Source reference.} Original lemma 05CJ: sources/stacks-chow/tags.txt, source lines 6922--6922. Cross-chapter standard prerequisite invoked once at proof lines 1877-1878 for the reduction step from N -> N' to N -> Sigma\textasciicircum{}\{-1\}N'.\par\medskip
\ifdefined\hypertarget\hypertarget{sourcepointer-4}{}\fi
\noindent\textit{Source reference.} Original lemma 00HQ: sources/stacks-chow/tags.txt, source lines 652--652. Cross-chapter standard prerequisite invoked once at proof lines 1887-1889 for faithful flatness of a ring map whose induced map on spectra hits all maximal ideals of a semi-local ring.\par\medskip
\ifdefined\hypertarget\hypertarget{sourcepointer-5}{}\fi
\noindent\textit{Source reference.} Original lemma 05CK: sources/stacks-chow/tags.txt, source lines 6923--6923. Cross-chapter standard prerequisite invoked once at proof lines 1894-1895 for universal injectivity of the module map induced by a faithfully flat ring map.\par\medskip
\ifdefined\hypertarget\hypertarget{sourcepointer-6}{}\fi
\noindent\textit{Source reference.} Original lemma 05CH: sources/stacks-chow/tags.txt, source lines 6920--6920. Cross-chapter standard prerequisite invoked together with 05CK at proof lines 1895-1896 to identify the localised map with N tensor\_S Sigma\textasciicircum{}\{-1\}S -> N tensor\_S Sigma\textasciicircum{}\{-1\}S'.\par\medskip
\ifdefined\hypertarget\hypertarget{sourcepointer-7}{}\fi
\noindent\textit{Source reference.} Original lemma 05CI: sources/stacks-chow/tags.txt, source lines 6921--6921. Cross-chapter standard prerequisite invoked once at proof lines 1900-1902 for closure of universal injectivity under composition; the final conclusion for N -> N' is drawn in the counted span.\par\medskip
\begin{lemma}
\label{lemma-base-change-universally-flat-local}
Let $R \to S$ be a ring map.
Let $N$ be an $S$-module.
Let $S \to S'$ be a ring map.
Assume
\begin{enumerate}
\item $R \to S$ is a local homomorphism of local rings
\item $S$ is essentially of finite presentation over $R$,
\item $N$ is of finite presentation over $S$,
\item $N$ is flat over $R$,
\item $S \to S'$ is flat, and
\item the image of $\Spec(S') \to \Spec(S)$ contains
all primes $\mathfrak q$ of $S$ lying over $\mathfrak m_R$
such that $\mathfrak q$ is an associated prime of $N/\mathfrak m_R N$.
\end{enumerate}
Then $N \to N \otimes_S S'$ is $R$-universally injective.
\end{lemma}

\begin{proof}
Set $N' = N \otimes_R S'$. Consider the commutative diagram
$$
\xymatrix{
N \ar[d] \ar[r] & N' \ar[d] \\
\Sigma^{-1}N \ar[r] & \Sigma^{-1}N'
}
$$
where $\Sigma \subset S$ is the set of elements which are not a
zerodivisor on $N/\mathfrak m_R N$. If we can show that the map
$N \to \Sigma^{-1}N'$ is universally injective, then $N \to N'$
is too (see
Algebra, Lemma \ref{algebra-lemma-universally-injective-permanence}).

\medskip\noindent
By
Lemma \ref{lemma-homothety-spectrum}
the ring $\Sigma^{-1}S$ is a semi-local ring whose maximal ideals
correspond to associated primes of $N/\mathfrak m_R N$.
Hence the image of
$\Spec(\Sigma^{-1}S') \to \Spec(\Sigma^{-1}S)$
contains all these maximal ideals by assumption. By
Algebra, Lemma \ref{algebra-lemma-ff-rings}
the ring map $\Sigma^{-1}S \to \Sigma^{-1}S'$ is faithfully flat.
Hence $\Sigma^{-1}N \to \Sigma^{-1}N'$, which is the map
$$
N \otimes_S \Sigma^{-1}S \longrightarrow N \otimes_S \Sigma^{-1}S'
$$
is universally injective, see
Algebra, Lemmas \ref{algebra-lemma-faithfully-flat-universally-injective} and
\ref{algebra-lemma-universally-injective-tensor}.
Finally, we apply
Lemma \ref{lemma-homothety-universally-injective}
to see that $N \to \Sigma^{-1}N$ is universally injective.
As the composition of universally injective module maps is universally
injective (see
Algebra, Lemma \ref{algebra-lemma-composition-universally-injective})
we conclude that $N \to \Sigma^{-1}N'$ is universally injective and we win.
\end{proof}

\end{document}
